% macros.tex
% ==================================================
% Helpers shared by every document.
% Load AFTER the application form, because these read its fields.
% ==================================================

%==================================================
% Draft mode
%   `make ... DRAFT=1` defines \draftmode. In draft mode every piece
%   of text that came from the application form is printed in color2,
%   so you can proofread exactly what was tailored, and a watermark
%   names the form (and mode) on every page.
%==================================================
\ifdefined\draftmode
  \newcommand{\draftcolor}[1]{{\color{color2}#1}}
  \usepackage{draftwatermark}
  \DraftwatermarkOptions{
    text  = {\sffamily DRAFT \textperiodcentered\ \letterfile
             \ifdefined\covermode\ \textperiodcentered\ \covermode\fi},
    color = color5, fontsize = 1.6cm, angle = 45}
\else
  \newcommand{\draftcolor}[1]{#1}
\fi

%==================================================
% \ifempty{\field}{if empty}{if filled}
%   The one test everything else is built on. A field counts as
%   empty if it is defined as {} or not defined at all, so you can
%   leave a line blank OR delete it from the form.
%==================================================
\newcommand{\ifempty}[3]{\ifdefvoid{#1}{#2}{#3}}

%==================================================
% Field printers
%   Use these inside paragraphs instead of the raw fields. Each one
%   prints its field (in color2 when drafting) or a sensible fallback
%   when the field is empty, so a letter never says "at ."
%==================================================
% Position: "the Data Scientist position" / "the advertised position"
\newcommand{\theposition}{\ifempty{\position}{advertised}{\draftcolor{\position}}}

% Employer: the company if set, otherwise the university,
% otherwise a neutral phrase. One paragraph serves both worlds.
% \employerarticle adds "the" in sentences ("at the University of
% Westacre") while the address block keeps the bare name.
\newcommand{\theemployer}{%
  \ifempty{\company}%
    {\ifempty{\uni}{your organization}{\draftcolor{\maybearticle\uni}}}%
    {\draftcolor{\maybearticle\company}}}
\newcommand{\maybearticle}{\ifempty{\employerarticle}{}{\employerarticle\ }}

% Short name for second mentions ("Westacre"); falls back to \theemployer
\newcommand{\theshortname}{\ifempty{\shortname}{\theemployer}{\draftcolor{\shortname}}}

%==================================================
% Optional-phrase macros
%   Each prints a whole phrase, leading space included, only when its
%   field is filled. Write them straight after a word, no space:
%     the \theposition\ position\maybeteam\ at \theemployer\maybejobref.
%   Empty team + empty job ref  ->  "the Analyst position at Spook Street."
%==================================================
\newcommand{\maybeteam}{\ifempty{\team}{}{ on the \draftcolor{\team} team}}
\newcommand{\maybedept}{\ifempty{\department}{}{ in the \draftcolor{\department}}}
\newcommand{\maybeschool}{\ifempty{\school}{}{ in the \draftcolor{\school}}}
\newcommand{\maybeproject}{\ifempty{\project}{}{ on the \draftcolor{\project} project}}
\newcommand{\maybespecialization}{\ifempty{\specialization}{}{ with a focus on \draftcolor{\specialization}}}
\newcommand{\maybejobref}{\ifempty{\jobref}{}{ (ref.\ \draftcolor{\jobref})}}
\newcommand{\maybeappsentence}{\ifempty{\appsentence}{}{ \draftcolor{\appsentence}}}

% Print a field (in draft color) only if it is filled.
\newcommand{\printifset}[1]{\ifempty{#1}{}{\draftcolor{#1}}}

%==================================================
% \coverpara{individualizedX}{file}
%   The paragraph override. If the form defines \individualizedX with
%   text, that text is used; if it is empty, the generic paragraph
%   modes/<mode>/<file>.tex is used instead.
%==================================================
\newcommand{\coverpara}[2]{%
  \ifcsvoid{#1}%
    {\modeinput{#2}}%
    {\draftcolor{\csname #1\endcsname}}%
  \par}

%==================================================
% "Prepared for" line under statement titles:
%   Position | Department | Employer
% Each part prints only if filled, so a fully generic copy (blank
% form) gets no line at all.
%==================================================
\newcommand{\prepitem}[1]{%
  \ifempty{#1}{}{\iffirstcontact\firstcontactfalse\else\contactsep\fi\mbox{\draftcolor{#1}}}}
\newcommand{\preparedfor}{%
  \begin{center}\vspace{-6pt}{\sffamily\small\color{color3}\firstcontacttrue
    \prepitem{\position}\prepitem{\department}%
    \ifempty{\company}{\prepitem{\uni}}{\prepitem{\company}}}\end{center}\vspace{-2pt}}
